% Part: proof-theory
% Chapter: normalization
% Section: reductions
%READERNOTE{SOL6-A305: न्यूनीकरणात सुरक्षित खुणा व आयगेन स्थिरांक वापरायचे. नव्या उघड्या गृहीतकांचा संच मूळच्या संचात समाविष्ट असतो; तो नेहमी तंतोतंत समान असेल किंवा खूण x पूर्णपणे नाहीशी होईल असा दावा नाही.}
%READERNOTE{SOL6-A309: नव्या परिच्छेदातील सिद्धता-नावे योग्य ग्रीक मुद्रण-संकेताने पुनर्स्थापित केली आहेत.}

\documentclass[../../../include/open-logic-section]{subfiles}

\begin{document}

\olfileid{pt}{nor}{red}

\olsection{न्यूनीकरण रूपांतरणे}

~\Log{N1i} मधील
!!a{proof} मध्ये एखाद्या
कटाची लांबी~$1$ असेल,
तर तो एका !!{formula}
आस्थितीचा खंड असतो.
तीच !!{formula} आस्थिती
!!a{introduction} अनुमानाचा किंवा~\FalseInt
चा निष्कर्ष आणि
!!a{elimination} अनुमानाचे
प्रमुख आधारविधान असते.
सामान्यरूपीकरणाच्या
सिद्धतेत अशा कटांचे
“न्यूनीकरण” करतात:
अशा !!{introduction}/!!{elimination}
जोडीत संपणारी उप-!!{proof}
त्या दोन्ही अनुमानांशिवाय
असलेल्या नव्या !!{proof}
ने बदलतात. याने नवे
कट निर्माण होऊ शकतात.
पण त्या नव्या कटांची
कोटी न्यूनीकृत कटाच्या
कोटीपेक्षा कमी असते.
त्यामुळे नव्या !!{proof}
ची कट-लांबी कमी होते
(कारण लांबी~$1$
असलेला महत्तम कट
निघून गेला), किंवा
!!{proof} ची कट-कोटी
कमी होते (संबंधित
महत्तम कट एकटाच
महत्तम कट असल्यास).

उदाहरणार्थ, संबंधित
अनुमाने~\Intro{\lif}
आणि~\Elim{\lif} असतील,
$!A \ident !B \lif !C$
असेल, आणि~$\delta_1$
ही उप-!!{proof}
पुढीलप्रमाणे संपत असेल:
\[
\AxiomC{$\Discharge{!B}{x}$}
\RightLabel{$\delta_2$}
\DeduceC{$!C$}
\DischargeRule{\Intro{\lif}}{x}
\UnaryInfC{$!B \lif !C$}
\AxiomC{}
\RightLabel{$\delta_3$}
\DeduceC{$!B$}
\RightLabel{\Elim{\lif}}
\BinaryInfC{$!C$}
\DisplayProof
\]
हा कट काढण्यासाठी
~$\delta$ मधील
उपसिद्धता~$\delta_1$
पुढील रचनेने बदलतो:
\[
\AxiomC{}
\RightLabel{$\delta_3$}
\DeduceC{$!B$}
\RightLabel{$\Subst{\delta_2}{\delta_3}{!B^x}$}
\DeduceC{$!C$}
\DisplayProof
\]
~$\Subst{\delta_2}{\delta_3}{!B^x}$
म्हणजे~$\delta_2$
मधील~$x$ खूण असलेल्या
प्रत्येक गृहीतक~$!B$
च्या जागी संपूर्ण
!!{proof}~$\delta_3$
ठेवून मिळणारी !!{proof}.
कलम-संयोजनापूर्वी $\delta_2$ च्या बद्ध मुक्ती-खुणा
आणि त्यांना बांधलेली गृहीतके एकत्र नवी नावे देऊन
$\delta_3$ च्या उघड्या खुणांपासून वेगळी ठेवायची.
त्यानंतर नव्या सिद्धतेची उघडी गृहीतके म्हणजे
$\delta_2$ मधील बसवलेल्या $!B^x$ आस्थितींव्यतिरिक्त
उरलेली उघडी गृहीतके, आणि प्रत्यक्ष बसवलेल्या
$\delta_3$ च्या प्रतींची उघडी गृहीतके.
$\delta_3$ मध्ये स्वतः $!B^x$ उघडे असू शकते;
म्हणून नव्या सिद्धतेत खूण $x$ अजिबात उरत नाही
असे म्हणता येत नाही. तसेच मुक्ती निष्फळ असेल किंवा
एखादी शाखा काढली गेली असेल, तर मूळची काही
उघडी गृहीतके गळू शकतात. आवश्यक निष्कर्ष असा की
नवी उघडी गृहीतके मूळ $\delta_1$ च्या उघड्या
गृहीतकांच्या पलीकडे जात नाहीत.
$\delta_2$ मधील मुक्ती-नियम आता आत बसवलेल्या
उघड्या गृहीतकांना अनवधानाने मुक्त करत नाहीत.
आणि~$\delta_3$
मधील गृहीतके मुक्त
करणारी अनुमाने अनेक
प्रतींमध्ये येऊ शकतात;
त्या प्रती
~$\Subst{\delta_2}{\delta_3}{!B^x}$
मधील संबंधित गृहीतके
मुक्त करतात.
सर्व न्यूनीकरण नमुन्यांत हीच सुरक्षित कलम-संयोजनाची
अट वापरायची. अनेक प्रतींतील आयगेन स्थिरांक आणि
त्यांच्याशी संबंधित बद्ध खुणा स्वतंत्र नवी नावे देऊन
नियमितता जपायची. संख्यापकांच्या नमुन्यांतील $t$ हे
घोषित बंद पदे वापरण्याच्या संकेतानुसार बंद पद आहे.

~$\delta_2$ किंवा
$\delta_3$ यांच्या
आतमध्ये पूर्णपणे असलेला
कट, किंवा~$\delta$
मध्ये~$\delta_1$
बाहेर पूर्णपणे असलेला
कट, बदलत नाही.
त्यात तीच !!{formula}s
आणि तीच अनुमाने
असतात; त्यामुळे त्याची
कोटी आणि लांबी
~$\delta$ मधील
अनुरूप कटासारखीच
राहते. आपण लांबी~$1$
असलेला महत्तम कट
काढला आहे; म्हणून
कट-लांबी कमी झाली,
किंवा~$\delta$
ची कट-लांबी~$1$
असून हा एकटाच महत्तम
कट होता, तर कट-कोटी
कमी झाली.

तथापि, दोन गोष्टी
घडू शकतात:
\begin{enumerate}
  \item सर्व~$!B^x$
  गृहीतकांच्या जागी
  $\delta_3$ ठेवल्याने
  ~$\delta_3$ मधील
  सर्व कटांच्या प्रती
  वाढतात. त्यांपैकी
  काही महत्तम कट
  असतील, तर नव्या
  !!{proof} ची कट-लांबी
  जुन्या !!{proof}
  पेक्षा अधिक होऊ
  शकते. हे होणार
  नाही याची खात्री
  हवी. म्हणून न्यूनीकरण
  रूपांतरणे फक्त
  \emph{सर्वात उजव्या}
  महत्तम कटांनाच
  लावतो: न्यूनीकृत
  होत असलेल्या कटाच्या
  उजवीकडे~$\delta$
  मधील कोणत्याही
  शाखेवर दुसरा महत्तम
  कट नसतो. ~$\delta_3$
  मध्ये महत्तम कट असता, तर
  तो येथे न्यूनीकृत
  होणाऱ्या कटाच्या
  उजवीकडे असता;
  म्हणून हा सर्वात
  उजवा कट नसता.
  सर्वात उजवे कट
  निवडल्याने !!{proof}
  ची कट-लांबी किंवा
  अगदी कट-कोटीही
  कमी होते.
  \item ~$\delta_2$
  मधील एखादे गृहीतक~$!B^x$
  !!a{elimination} अनुमानाचे
  प्रमुख आधारविधान
  असेल, आणि~$\delta_3$
  चे अंतिम अनुमान
  \Elim{\lor}, \Elim{\lexists}
  किंवा !!a{introduction}
  अनुमान असेल, तर
  त्या गृहीतकावर
  ~$\delta_3$ चे
  कलम-संयोजन केल्यामुळे
  नवा कट तयार होतो.
  ~$\delta_2$ मध्ये
  जे केवळ गृहीतक
  होते, ते आता लांबी~$1$
  असलेला कट
  (!!a{introduction}
  अनुमानाचा निष्कर्ष आणि
  !!a{elimination} अनुमानाचे
  प्रमुख आधारविधान)
  किंवा कट-खंडाची
  शेवटची !!{formula}
  आस्थिती
  (!!a{elimination}
  अनुमानाचे प्रमुख
  आधारविधान आणि
  \Elim{\lor}
  किंवा~\Elim{\lexists}
  चा निष्कर्ष)
  होते. $B^x$ हे
  \Elim{\lor}
  किंवा~\Elim{\lexists}
  चे गौण आधारविधान
  असेल, तर ते आधीच
  एखाद्या खंडाची,
  कदाचित कट-खंडाचीही,
  पहिली !!{formula}
  आस्थिती होती.
  नव्या !!{proof}
  मध्ये तो खंड
  आता अधिक लांब
  होऊ शकतो. पण
  या नव्या कटाचे
  किंवा विद्यमान
  कट-खंडाचे !!{formula}
  हे~$!B$ आहे, आणि
  $\depth{!B} < \depth{!B \lif !C}$.
  म्हणून नवे कट
  आले किंवा काही
  कटांची लांबी
  वाढली, तरी त्यांपैकी
  कोणताही महत्तम कट
  नसतो. अशा रीतीने
  कट-कोटी कमी होते,
  किंवा समान कोटीचे
  कट उरले असतील
  तर कट-लांबी कमी
  होते.
\end{enumerate}

लांबी~$1$ असलेला
कट कमी करण्याच्या
प्रक्रियेला
\emph{न्यूनीकरण रूपांतरण}
(किंवा
\emph{वळसा रूपांतरण})
म्हणतात. काही
नियम-जोड्यांचे नमुने
~\olref{tab:reduction-conversions}
मध्ये दिले आहेत.
(उरलेले सरावासाठी
सोडले आहेत.)

\begin{table}
\begin{multline*}
\bottomAlignProof
\AxiomC{$\Discharge{!B}{x}$}
\RightLabel{$\delta_2$}
\shortDeduce
\DeduceC{$!C$}
\DischargeRule{\Intro{\lif}}{x}
\UnaryInfC{$!B \lif !C$}
\AxiomC{}
\RightLabel{$\delta_3$}
\shortDeduce
\DeduceC{$!B$}
\RightLabel{\Elim{\lif}}
\BinaryInfC{$!C$}
\DisplayProof
\quad \longrightarrow\quad
\shoveright{\bottomAlignProof
\AxiomC{}
\RightLabel{$\delta_3$}
\shortDeduce
\DeduceC{$!B$}
\RightLabel{$\Subst{\delta_2}{\delta_3}{!B^x}$}
\shortDeduce
\DeduceC{$!C$}
\DisplayProof}
\\[4ex]
\shoveleft{\bottomAlignProof
\AxiomC{}
\RightLabel{$\delta_2$}
\shortDeduce
\DeduceC{$!B$}
\AxiomC{}
\RightLabel{$\delta_3$}
\shortDeduce
\DeduceC{$!C$}
\RightLabel{\Intro{\land}}
\BinaryInfC{$!B \land !C$}
\RightLabel{\Elim{\land}[1]}
\UnaryInfC{$!B$}
\DisplayProof
\quad \longrightarrow\quad
\bottomAlignProof
\AxiomC{}
\RightLabel{$\delta_2$}
\shortDeduce
\DeduceC{$!B$}
\DisplayProof}\qquad
\shoveright{\bottomAlignProof
\AxiomC{}
\RightLabel{$\delta_2$}
\shortDeduce
\DeduceC{$!B$}
\AxiomC{}
\RightLabel{$\delta_3$}
\shortDeduce
\DeduceC{$!C$}
\RightLabel{\Intro{\land}}
\BinaryInfC{$!B \land !C$}
\RightLabel{\Elim{\land}[2]}
\UnaryInfC{$!C$}
\DisplayProof
\quad \longrightarrow\quad
\bottomAlignProof
\AxiomC{}
\RightLabel{$\delta_3$}
\shortDeduce
\DeduceC{$!C$}
\DisplayProof}
\\[4ex]
\bottomAlignProof
\AxiomC{}
\RightLabel{$\delta_2$}
\shortDeduce
\DeduceC{$!B$}
\RightLabel{\Intro{\lor}[1]}
\UnaryInfC{$!B \lor !C$}
\AxiomC{$\Discharge{!B}{x}$}
\RightLabel{$\delta_3$}
\shortDeduce
\DeduceC{$!D$}
\AxiomC{$\Discharge{!C}{y}$}
\RightLabel{$\delta_4$}
\shortDeduce
\DeduceC{$!D$}
\DischargeRule{\Elim{\lor}}{x\,y}
\TrinaryInfC{$!D$}
\DisplayProof
\quad \longrightarrow\quad
\shoveright{\bottomAlignProof
\AxiomC{}
\RightLabel{$\delta_2$}
\shortDeduce
\DeduceC{$!B$}
\RightLabel{$\Subst{\delta_3}{\delta_2}{!B^x}$}
\shortDeduce
\DeduceC{$!D$}
\DisplayProof}
\\[2ex]
\shoveleft{\bottomAlignProof
\AxiomC{}
\RightLabel{$\delta_2$}
\shortDeduce
\DeduceC{$!B(c)$}
\RightLabel{\Intro{\lforall}}
\UnaryInfC{$\lforall[x][!B(x)]$}
\RightLabel{\Elim{\lforall}}
\UnaryInfC{$!B(t)$}
\DisplayProof
\quad \longrightarrow\quad
\bottomAlignProof
\AxiomC{}
\RightLabel{$\Subst{\delta_2}{t}{c}$}
\shortDeduce
\DeduceC{$!B(t)$}
\DisplayProof}
\\[4ex]
\shoveright{\bottomAlignProof
\AxiomC{}
\RightLabel{$\delta_2$}
\shortDeduce
\DeduceC{$!B(t)$}
\RightLabel{\Intro{\lexists}}
\UnaryInfC{$\lexists[x][!B(x)]$}
\AxiomC{$\Discharge{!B(c)}{x}$}
\RightLabel{$\delta_3$}
\shortDeduce
\DeduceC{$!D$}
\DischargeRule{\Elim{\lexists}}{x}
\BinaryInfC{$!D$}
\DisplayProof
\quad \longrightarrow\quad
\bottomAlignProof
\AxiomC{}
\RightLabel{$\delta_2$}
\shortDeduce
\DeduceC{$!B(t)$}
\RightLabel{$\Subst{\Subst{\delta_3}{t}{c}}{\delta_2}{!B(t)^x}$}
\shortDeduce
\DeduceC{$!D$}
\DisplayProof}
\end{multline*}
\caption{\Log{N1i} साठी काही न्यूनीकरण रूपांतरणे.}
\ollabel{tab:reduction-conversions}
\end{table}

\begin{prob}
\Intro{\lnot} नंतर
\Elim{\lnot} येत
असेल, तर न्यूनीकरण
रूपांतरण द्या.
\end{prob}

आणखी एक प्रसंग
पाहावा लागतो. कटाच्या
व्याख्येत असा प्रसंग
येऊ शकतो की कटाची
लांबी~$1$ आहे,
आणि~$!A$ हे
\FalseInt{} चा निष्कर्ष
तसेच !!a{elimination}
अनुमानाचे प्रमुख
आधारविधान आहे.
$!A$ हा
!!a{introduction}
नियमाचा निष्कर्ष
असताना करतो तसेच
येथेही करू. उदाहरणार्थ,
पुढील बदल करा:
\[
\bottomAlignProof
\AxiomC{}
\RightLabel{$\delta_2$}
\DeduceC{$\lfalse$}
\UnaryInfC{$!B \lif !C$}
\AxiomC{}
\RightLabel{$\delta_3$}
\DeduceC{$!B$}
\RightLabel{\Elim{\lif}}
\BinaryInfC{$!C$}
\DisplayProof
\quad\text{ऐवजी}\quad
\bottomAlignProof
\AxiomC{}
\RightLabel{$\delta_2$}
\DeduceC{$\lfalse$}
\UnaryInfC{$!C$}
\DisplayProof
\]
~$!A \ident (!B \lor
!C)$ किंवा
$!A \ident \lexists[x][!B(x)]$
असेल, तर थोडी काळजी
हवी. उदाहरणार्थ,
पुढील रचनेच्या जागी
\[
\bottomAlignProof
\AxiomC{}
\RightLabel{$\delta_2$}
\DeduceC{$\lfalse$}
\RightLabel{\FalseInt}
\UnaryInfC{$!B \lor !C$}
\AxiomC{$\Discharge{!B}{x}$}
\RightLabel{$\delta_3$}
\DeduceC{$!D$}
\AxiomC{$\Discharge{!C}{y}$}
\RightLabel{$\delta_4$}
\DeduceC{$!D$}
\DischargeRule{\Elim{\lor}}{x\,y}
\TrinaryInfC{$!D$}
\DisplayProof
\quad\text{ऐवजी}\quad
\bottomAlignProof
\AxiomC{}
\RightLabel{$\delta_2$}
\DeduceC{$\lfalse$}
\RightLabel{\FalseInt}
\UnaryInfC{$!D$}
\DisplayProof
\]
असा बदल केला,
तर कट-!!{formula}~$!D$
असलेला नवा कट
निर्माण होऊ शकतो.
पण
$\depth{!D} < \depth{!B \lor !C}$
असेलच याची हमी नाही!{}
म्हणून~\Intro{\lor}/\Elim{\lor}
न्यूनीकरणात जसे
केले, तसेच येथेही
करावे आणि पुढील
रचनेने बदलावे:
\[
\bottomAlignProof
\AxiomC{}
\RightLabel{$\delta_2$}
\DeduceC{$\lfalse$}
\RightLabel{\FalseInt}
\UnaryInfC{$!B$}
\RightLabel{$\delta_3$}
\DeduceC{$!D$}
\DisplayProof
\quad\text{किंवा}\quad
\bottomAlignProof
\AxiomC{}
\RightLabel{$\delta_2$}
\DeduceC{$\lfalse$}
\RightLabel{\FalseInt}
\UnaryInfC{$!C$}
\RightLabel{$\delta_4$}
\DeduceC{$!D$}
\DisplayProof
\]

\end{document}
